\section{Задание}

\textbf{Цель работы:} ознакомление с базовыми возможностями Verilog HDL на примере создания цифрового устройства из нескольких модулей, реализующих арифметические операции, с передачей результатов вычислений по сигналу синхронизации между модулями.

В процессе выполнения лабораторной работы необходимо:
\begin{itemize}
    \item Разработать два варианта реализации БОЭ на языке Verilog в соответствии с вариантом.
    \begin{itemize}
        \item Используя только assign конструкции
        \item Используя always-блок
    \end{itemize}
    \item Разработать расширяемый БОЭ в виде параметризируемого модуля с возможностью настраивать ширину входных аргументов. Использовать уже разработанный БОЭ и generate-блок.
    \item Разработать последовательный элемент для сохранения результата расширяемого комбинационного элемента в регистр по тактовому сигналу \texttt{clk}. Поддержать возможность сброса значения результата по сигналу \texttt{rst}.
    \item Протестировать разработанный элемент в симуляционном окружении.
    \item Составить отчет по результатам выполнения работы.
\end{itemize}

Вариант работы определяется последней цифрой идентификатора ИСУ студента. В случае, если последняя цифра выходит за пределы области значений вариантов, используется следующая цифра идентификатора ИСУ, расположенная слева от последней. Итерации повторяются до тех пор, пока не будет найдено соответствие.

Таким образом, в моём случае вариант работы равен:
\[
    \texttt{get\_variant}(466449) = 4.
\]

Этому варианту соответствует базовый операционный элемент \textbf{дешифратора 1 в 2}.